% ===================== 第 4 章（中） 翻译真题 =====================
\section{真题分类汇编：Translation}

\begin{tipbox}[title=本章说明 / Note]
汉译英部分：\textbf{题目原文为中文段落（180--200 汉字），保留中文是为了保留完整题型}；参考译文（Reference version）为本指南自编译文，供对照思路使用，并非官方答案。话题索引整理自公开真题资料。
\end{tipbox}

\subsection{Past-exam Topic Index (2013--2025)}

\begin{center}
\small
\begin{longtable}{p{2.3cm}p{6.2cm}p{6.2cm}}
\toprule
\textbf{Date} & \textbf{Topic} & \textbf{Category} \\
\midrule
\endfirsthead
\toprule
\textbf{Date} & \textbf{Topic} & \textbf{Category} \\
\midrule
\endhead
2013.12 & The Silk Road; Chinese food; tea & Culture / History \\
\midrule
2014.06 & China's economic development; education & Economy / Society \\
\midrule
2015.06 & Chinese diet (rice and wheat); water conservancy & Life / Geography \\
\midrule
2016.12 & Mount Tai / Mount Huang / the Three Gorges & Geography \\
\midrule
2017.06 & The Yangtze River / the Yellow River / the Pearl River & Geography \\
\midrule
2018.06 & Public libraries; museums; public transport & Society \\
\midrule
2019.06 & Paper-cutting; Chinese painting; shadow puppetry & Intangible culture \\
\midrule
2019.12 & Plum blossom / peony / lotus; flowers in Chinese culture & Culture \\
\midrule
2020.09 & Water Margin; Dream of the Red Chamber; Journey to the West & Four great classical novels \\
\midrule
2020.12 & Daxing International Airport; high-speed rail & Infrastructure \\
\midrule
2021.06 & Hainan Free Trade Port; Yan'an; rural development & Society / Economy \\
\midrule
2022.06 & Zhaozhou Bridge; Nanjing Yangtze River Bridge; Lugou Bridge & Geography / History \\
\midrule
2023.06 & Rural revitalization; government policies & Society / Economy \\
\midrule
2024.06 & Traditional Chinese opera (singing, reciting, acting and acrobatic fighting) & Intangible culture \\
\midrule
2025.06 & Tiangong Space Station; space science & Science / Technology \\
\midrule
2025.12 & Respecting the elderly; frugality; the spirit of self-reliance & Traditional virtues \\
\bottomrule
\end{longtable}
\end{center}

\begin{notebox}[title=指数式话题分布（近十年）]
Culture and history $\approx$ 40\% \quad Society and economy $\approx$ 35\% \quad Geography and daily life $\approx$ 25\%. The current trend is ``traditional culture + contemporary hot topics'' (e.g.\ intangible heritage + short videos, craftsmanship + digital protection).
\end{notebox}

\subsection{Selected Real Items with Reference Versions}

\subsubsection*{Item 1 (2025.12, Set 1): Respecting the Elderly}

\textbf{Chinese source (真题原文):}

\begin{quote}
尊老是中华民族的传统美德，深深植根于中国人的思想和行为中，是人们普遍遵守的行为规范和社会准则。在当今的中国，这种美德得到广泛传承。社会各界积极营造敬老助老的社会氛围，为老年人提供便利服务已成为社会共识，例如社区专门开设长者食堂，公共场所配置优先座位。政府还出台了一系列政策，为老年人的权益提供有力保障。尊老是社会和谐与发展的重要基石，有助于培育良好的社会氛围，推动整个社会文明进步。
\end{quote}

\textbf{Reference version:}

\begin{engbox}
Respecting the elderly is a traditional virtue of the Chinese nation, deeply rooted in the thoughts and behaviour of the Chinese people and widely observed as a code of conduct and a social norm. In today's China, this virtue is being carried forward on a broad scale. All sectors of society are working to create an atmosphere of respecting and helping the elderly, and providing convenient services for them has become a social consensus: communities have opened canteens especially for senior citizens, and priority seats are available in public places. The government has also introduced a series of policies that effectively protect the rights and interests of the elderly. Respecting the elderly is an important cornerstone of social harmony and development, helping to foster a sound social atmosphere and advance the civilization of the whole society.
\end{engbox}

\subsubsection*{Item 2 (2025.12, Set 2): Frugality}

\textbf{Chinese source (真题原文):}

\begin{quote}
节俭是中华民族的传统美德。在古代中国，生产力低，我们的祖先深知劳动成果来之不易，因此遵循节俭的理念。在物质极其丰富的今天，中国人仍然坚持节俭的生活方式，这体现了理性消费和尊重劳动的哲学。近年来，中国政府不断加大节约型社会建设力度，倡导“清盘运动”以减少食物浪费，倡导简单低碳的生活方式，推进资源循环利用，反对过度消费。发展以节约为导向的社会，已切实推动了中国经济的高质量发展。
\end{quote}

\textbf{Reference version:}

\begin{engbox}
Frugality is a traditional virtue of the Chinese nation. In ancient China, productivity was low, and our ancestors knew well that the fruits of labour were hard-won, so they followed the idea of thrift. In today's world of abundant material wealth, the Chinese people still stick to a frugal way of life, which reflects a philosophy of rational consumption and respect for labour. In recent years, the Chinese government has stepped up efforts to build a conservation-oriented society: it advocates the ``Clean Plate Campaign'' to reduce food waste, promotes a simple, low-carbon lifestyle, advances resource recycling, and opposes excessive consumption. The development of such a conservation-oriented society has effectively boosted the high-quality growth of China's economy.
\end{engbox}

\subsubsection*{Item 3 (2025.12, Set 3): The Spirit of Self-reliance}

\textbf{Chinese source (真题原文):}

\begin{quote}
自力更生的精神深深植根于中国文化之中，这一精神已成为国家生存与发展的关键。多年来，中国在自强和自力更生的道路上持续加大对科研领域的投入，并在信息技术等领域取得了新的突破。全国范围内的5G网络覆盖、远程医疗的发展以及嫦娥探月工程，都展示了中国在航天与科技领域取得的显著进展。自力更生这一强大动力激励着中国人民为实现民族复兴而奋斗。
\end{quote}

\textbf{Reference version:}

\begin{engbox}
The spirit of self-reliance is deeply rooted in Chinese culture and has become crucial to the nation's survival and development. Over the years, on the path of self-reliance and independent development, China has continuously increased its investment in scientific research and made new breakthroughs in fields such as information technology. The nationwide coverage of 5G networks, the development of telemedicine, and the Chang'e lunar exploration programme all demonstrate China's remarkable progress in aerospace and technology. This powerful driving force of self-reliance inspires the Chinese people to strive for national rejuvenation.
\end{engbox}

\subsubsection*{Item 4 (2025.06, Set 1): Tiangong Space Station}

\textbf{Chinese source (真题原文):}

\begin{quote}
天宫空间站是中国太空实验室，拥有110多立方米使用空间，可驻留三名宇航员，在距地球表面400--450公里的轨道上运行。天宫空间站已实施180多个科学研究与应用项目，涉及空间生命科学、太空医学、空间材料科学等领域。天宫空间站的研究成果在我国得到了广泛应用，产生了显著的经济效益。例如，太空育种创造的直接经济效益高达数千亿元。这不仅标志中国在航天技术上取得了巨大进步，也表明中国为全球的太空研究和应用做出了重大贡献。
\end{quote}

\textbf{Reference version:}

\begin{engbox}
The Tiangong Space Station is China's space laboratory. With more than 110 cubic metres of usable space, it can host three astronauts and operates in an orbit 400--450 kilometres above the Earth's surface. The station has carried out over 180 scientific research and application projects, covering fields such as space life sciences, space medicine and space materials science. Its research results have been widely applied in China and have produced remarkable economic benefits; space breeding, for example, has generated direct economic benefits worth hundreds of billions of yuan. This not only marks China's tremendous progress in aerospace technology but also shows that China has made significant contributions to global space research and applications.
\end{engbox}

\subsubsection*{Item 5 (2020.09, Set 1): Water Margin}

\textbf{Chinese source (真题原文):}

\begin{quote}
《水浒传》是中国文学四大经典小说之一。这部小说基于历史人物宋江及其伙伴反抗封建帝王的故事，数百年来一直深受中国读者的喜爱。毫不夸张地说，几乎每个中国人都熟悉小说中的一些主要人物。这部小说中的精彩故事在茶馆、戏剧舞台、广播电视、电影屏幕和无数家庭中反复讲述。事实上，这部小说的影响已经远远超出了国界，越来越多的外国读者也感到这部小说里的故事生动感人、趣味盎然。
\end{quote}

\textbf{Reference version:}

\begin{engbox}
Water Margin is one of the four great classical novels of Chinese literature. Based on the story of the historical figure Song Jiang and his companions rebelling against the feudal emperor, it has been popular among Chinese readers for centuries. It is no exaggeration to say that almost every Chinese is familiar with some of its main characters. The novel's wonderful stories have been retold again and again in teahouses, on the stage, on radio and television, on cinema screens and in countless homes. In fact, its influence has long extended far beyond national borders, and more and more foreign readers also find its stories vivid, moving and full of interest.
\end{engbox}

\subsection{Translation Tips (English note form)}

\begin{itemize}
  \item \textbf{Topic first.} Decide within ten seconds whether the passage is about culture, society or geography, then activate the matching word bank from Chapter 3.
  \item \textbf{Proper nouns matter.} Use standard spellings: \emph{The Analects of Confucius}, the Yangtze River, the Qinghai-Tibet Plateau, the 24 Solar Terms, intangible cultural heritage, the Reform and Opening-up policy.
  \item \textbf{Convert parts of speech.} Chinese verbs often become English nouns: 他很喜欢读书 $\to$ He has a great love for reading.
  \item \textbf{Prefer the passive when it fits.} 被誉为 / 被认为 / 被列入 $\to$ be praised as / be regarded as / be listed as.
  \item \textbf{Combine short clauses.} Chinese chains short sentences; English prefers one main clause plus participles or attributive clauses.
  \item \textbf{Do not translate word for word.} Four-character phrases usually need explanation, not literal rendering (丰富多彩 $\to$ rich and varied).
\end{itemize}
